\begin{proposition}[CJ from the active phase-event law]
Keep the theorem object
\[
B=(s,u,v,\lambda,f),\qquad \tau,\qquad \epsilon_4.
\]
Define the active phase
\[
\Theta:=q+s+v+\lambda\pmod m.
\]
Assume the active branch satisfies the phase-event law
\begin{align*}
\Theta=0 &\Longrightarrow \epsilon_4=\mathrm{wrap},\\
\Theta=1 &\Longrightarrow \epsilon_4=\mathrm{carry\_jump},\\
\Theta=2 &\Longrightarrow \epsilon_4=\mathrm{other}\iff c=1,\\
\Theta=3 &\Longrightarrow \epsilon_4=\mathrm{other}\iff \bigl(c=0\text{ and }s=2\bigr)\text{ or }\bigl(c=1\text{ and }s\neq 2\bigr),\\
4\le \Theta\le m-1 &\Longrightarrow \epsilon_4=\mathrm{flat}.
\end{align*}
Then the flat-corner law follows:
\[
\delta=1 \Longrightarrow \text{next event is wrap},
\]
\[
2\le \delta\le m-3 \Longrightarrow \text{next event is flat},
\]
\[
\delta=m-2 \Longrightarrow
\begin{cases}
\text{next event is other},& s=2,\\
\text{next event is flat},& s\neq 2.
\end{cases}
\]
Consequently, the carry\_jump reset theorem follows:
\[
R_{\mathrm{cj}}=
\begin{cases}
0,& s+v+\lambda\equiv 2\pmod m,\\
1,& s=1\text{ and }s+v+\lambda\not\equiv 2\pmod m,\\
m-2,& \text{otherwise}.
\end{cases}
\]
\end{proposition}

\begin{proof}[Proof sketch]
On flat states one has
\[
\delta\equiv -\Theta\pmod m.
\]
If \(\delta=1\), then \(\Theta=m-1\) and the successor has phase \(0\), hence is wrap.
If \(2\le \delta\le m-3\), then the successor phase lies in \(\{4,\dots,m-1\}\), hence is flat.
If \(\delta=m-2\), then the current flat state has \(\Theta=2\), so the phase-event law forces \(c=0\); its successor has \(\Theta=3\) with the same \(s\) and \(c=0\), hence is other iff \(s=2\).
This is exactly the flat-corner law.  The previously established carry\_jump reduction then yields the stated formula for \(R_{\mathrm{cj}}\).
\end{proof}
